\section{Electromagnetism}

The rule that streams and collides the tone already conserves a charge. When that conservation is read at the level of the lattice it behaves as a discrete Gauss law, and the field it sources turns out to be an emergent $U(1)$ gauge field with the structure physicists expect. What follows shows the conservation law, the gauge field built on top of it, the recovery of the Coulomb potential, and the gauge-invariant observables that pin down the field's physical content. None of this is added by hand. It is what the charge rule does once the bookkeeping is followed carefully.

\subsection{The Gauss law}

The base rule conserves charge \emph{locally}. The charge that flows into a cell over one beat equals the charge that flows out, cell by cell, with no global ledger needed. That statement is a discrete continuity equation, and a discrete continuity equation is a discrete Gauss law. The conserved current it defines is exactly the source term for the emergent $U(1)$ field. The conservation holds exactly and on every substrate tested, not approximately and not only in some limit, because it is built into the STREAM and COLLIDE moves rather than imposed afterward.

\begin{result}[Discrete Gauss law]
The base rule conserves charge per cell per beat. This local conservation is a discrete continuity equation whose conserved current sources an emergent $U(1)$ field. It holds exactly on every substrate examined \cite{pollard2026vibetest}.
\end{result}

\subsection{The emergent gauge field}

To probe the field directly one places a vortex gauge configuration on the lattice. Link phases wind around a plaquette so that the loop encloses a flux $\Phi$. The Wilson loop, the sum of the link phases taken around a closed path, then equals the enclosed flux. The value is gauge-invariant. Applying a random gauge transformation $A \to A + d\lambda$ leaves the Wilson loop unchanged, which is the signature of a genuine gauge theory rather than a coordinate artifact. The Aharonov-Bohm phase picked up by a charge that encircles the flux is exactly $\Phi$, independent of the shape of the loop and of the gauge. A loop that does not enclose the flux gives zero. The physical content of the field is therefore its gauge-invariant field strength, and the field is a $U(1)$ gauge theory in the ordinary sense, which is to say electromagnetism.

\begin{result}[Gauge-invariant $U(1)$ observables]
On a vortex configuration the Wilson loop equals the enclosed flux $\Phi$, is invariant under $A \to A + d\lambda$, and produces an Aharonov-Bohm phase $\Phi$ for any encircling loop and zero for a non-encircling one \cite{pollard2026vibetest}. The emergent field is a $U(1)$ gauge theory.
\end{result}

\subsection{The Coulomb potential}

The static sector of the emergent field is governed by the flat three-dimensional Laplacian. Its Green's function is the familiar $1/r$ Coulomb potential, recovered here by the same difference method used for the Newtonian potential on the flat cusp. The static $U(1)$ field around a point charge falls off as $1/r$, with no free constants beyond the overall scale. The electromagnetic and gravitational static potentials share this origin. Both are Green's functions of the cusp Laplacian.

\subsection{The non-abelian extension}

The Wilson-loop construction does not depend on the abelian case. Replacing the link phases with matrices gives a non-abelian gauge field, the same matrices that appear in the $\mathrm{SO}(10)$ analysis of the coin-plus-tone sector. The Wilson loop is again gauge-invariant, the field is curved, and the link variables no longer commute. The abelian electromagnetic field is the simplest member of a family that the same machinery produces, which is the expected relationship between $U(1)$ and the larger gauge groups of the Standard Model.


Electromagnetism here is emergent rather than postulated. The charge rule's local conservation forces a $U(1)$ gauge symmetry, and the field that symmetry organizes carries the Coulomb $1/r$ potential and the gauge-invariant Wilson loops one would demand of a real gauge theory. That closes the base-to-gauge step for the $U(1)$ sector.
